\section{静态快照防漂移契约与工程难点攻关}\label{sec:contracts-challenges}

\subsection{构建管线 SHA-256 指纹互锁与防漂移机制}

在多人同时迭代统计脚本、科学图稿与前端页面的几周里，最容易出纰漏的地方就是“上游重跑了某个 CSV，下游网页数据包却忘了同步更新”。为了把这种人为疏忽彻底拦在构建阶段，我们在 \texttt{scripts/system/build\_web\_app.py} 与 \texttt{check\_web\_snapshot.py} 中建立了一套严格的哈希互锁契约：

\begin{itemize}
  \item \textbf{图证同源强校验}：每次执行构建或运行 \texttt{python3 app/launch.py} 时，脚本都会重新计算 \texttt{code-evidence.json} 与压缩特征包 \texttt{code-features.json.gz} 的字节级 SHA-256 摘要，并与 \texttt{figure-metadata.json} 中记录的源指纹逐一比对；一旦发现图稿与数据包不属于同一次快照生成，构建器当场抛出异常终止打包，绝不允许新旧口径混装上线。
  \item \textbf{静态产物内容寻址防缓存}：所有拷贝至 \texttt{app/dist/} 运行目录的样式表与脚本均自动追加由文件内容计算出的 12 位版本查询串（如 \texttt{?v=64c2ad0fc74d}），同时对内置的四大前端可视化库核验许可证与包哈希，确保本地离线打开时不会读到陈旧的浏览器磁盘缓存。
  \item \textbf{科研资产保护白名单}：脚本 \texttt{check\_reconstruction.py} 对仓库内 37 幅核心主题 SVG、12 份数据源清单及总计 \textbf{1,100 项已验收科研文件指纹}建立只读保护监控，确保美化界面交互时不会无意篡改任何一条实验曲线的坐标数值。
\end{itemize}
